\subsection{Characterwise reciprocity does not supply admissibility}

\begin{claim}{D.29}{Proved; a limitation of the reciprocity route}\label{clm:D29}
There are elements $A_p\in\Q_p^\times$ at every finite
rational prime and $A_\infty=1$ such that the following
statements hold.
\begin{enumerate}
\item The values $x_v=2+A_v$ are ordinary points of
$\mathbb P^1\setminus\{0,1,2,\infty\}$ over $\Q_v$.
\item For every $n>1$, the tuple
$(\kappa_n(A_v))_v$ lies outside the restricted product
$P^1(\Q,\mu_n)$. No global element of
$\widehat{\Q^\times}$ has these localizations.
\item For every number field $L\subset\overline\Q$, every
$n\ge1$, and every $\chi\in H^1(L,\Z/n)$, the sum
\[
 \sum_w\operatorname{inv}_w
       \bigl(\kappa_n(A_{L,w})\smile\chi_w\bigr)           \tag{D.29.1}
\]
has finitely many nonzero summands and equals zero. Here
$A_{L,w}$ is the image of $A_p$ for $w\mid p$ and is $1$
at every archimedean place. Nevertheless, this tuple lies
outside $P^1(L,\mu_n)$ for every $n>1$.
\end{enumerate}
The exceptional rational primes in the second assertion
have Dirichlet density zero. At each of them both original
tripod coordinates $x_p$ and $1-x_p$ are units. The data
are compatible with restriction to finite number fields and
with the usual norm-of-restriction identity.

These points are not asserted to come from a section, to
satisfy global completed Kummer conditions for $t$ and
$1-t$, or to satisfy the conditions on all accessible finite
\'etale neighbourhoods. The claim proves that characterwise
reciprocity, even after every finite constant extension,
does not replace the restricted-product hypothesis in D.27.
\end{claim}
\begin{proof}
Enumerate the rational primes as $p_1,p_2,\ldots$ in
increasing order. Enumerate finite Galois extensions
$E_j/\Q$ inside a fixed algebraic closure, with every such
extension eventually occurring. For a reproducible numerical
prefix choose
\[
 E_1=\Q,\qquad E_2=\Q(i),\qquad
 E_3=E_4=\Q(\zeta_{24}),\qquad E_5=\Q(\zeta_{120}),
\]
and then enumerate the splitting fields of all monic integer
polynomials. Repetition is harmless.

Set $P_0=M_0=q_0=1$. Given the preceding choices, define
\[
\begin{split}
 R_j&=1+\max\{v_{p_j}(P_{j-1}+1),
                         v_{p_j}(P_{j-1}+2)\},\\
 M_j&=\operatorname{lcm}
           \bigl(M_{j-1},(j!)^j,p_j^{R_j}\bigr).
\end{split}                                                \tag{D.29.2}
\]
Choose $q_j$ to be the least prime exceeding
$\max\{2^j,q_{j-1},p_j\}$ that splits completely in
\[
                    E_1\cdots E_j\,\Q(\mu_{M_j}),
\]
and put $P_j=P_{j-1}q_j$. The field displayed is finite
Galois over $\Q$. Chebotarev for the identity conjugacy
class supplies infinitely many such primes
\cite[Chapter V, Theorem 3.23, p.~164]{D_MilneCFT2020}.
Thus the recursion is defined. Each step involves finite
exact splitting tests; it does not require deciding an
infinite splitting condition. In particular
\[
 q_j\equiv1\pmod {M_j},\qquad q_j>2^j,                     \tag{D.29.3}
\]
and every sufficiently late $q_j$ splits completely in any
one prescribed finite Galois extension of $\Q$.

\emph{Local limits and actual point values.}
Fix a prime $p$. For every $r\ge1$, all sufficiently late
factors $q_j$ lie in $1+p^r\Z_p$, because
$(j!)^j\mid M_j$. Arbitrary finite products of these factors
remain in that subgroup. After a fixed finite initial
product, the sequence $(P_N)_N$ is therefore Cauchy and its
tail converges to a unit. Exactly one factor equals $p$ if
$p$ is among the $q_j$, and none does otherwise. Hence
\[
 A_p:=\lim_{N\to\infty}P_N\in\Q_p^\times,
 \qquad
 v_p(A_p)=
 \begin{cases}
  1,&p=q_j\text{ for some }j,\\
  0,&\text{otherwise}.
 \end{cases}                                               \tag{D.29.4}
\]
There is no claim of real convergence of the positive
products $P_N$. We separately set $A_\infty=1$; all $P_N$
have its real Kummer class for every modulus.

The extra exponent $R_j$ ensures that the local limits
avoid the two forbidden translated values. For $p=p_j$,
the divisibility $M_j\mid M_i$ for $i\ge j$ implies
\[
              A_p/P_{j-1}\in1+p^{R_j}\Z_p.
\]
For $c=1,2$ this gives
\[
 v_p(A_p-P_{j-1})\ge v_p(P_{j-1})+R_j
                         >v_p(P_{j-1}+c).
\]
Consequently
$v_p(A_p+c)=v_p(P_{j-1}+c)<\infty$.
In particular $A_p\ne-1,-2$. Since also $A_p\ne0$,
the point $x_p=2+A_p$ is finite and different from $0,1,2$.
The real choice $x_\infty=3$ has the same property.

\emph{Failure of the restricted product.}
For any $n>1$, infinitely many $q_j$ are prime to $n$.
At each of them $\kappa_n(A_{q_j})$ has valuation
$1\bmod n$. An unramified Kummer class at a place of
residue characteristic prime to $n$ has valuation zero
modulo $n$. Thus the tuple is ramified at infinitely many
places and lies outside $P^1(\Q,\mu_n)$.

Every rational number has nonzero valuation at only
finitely many primes. Its squareclass cannot have the
localizations in (D.29.4). The modulus-two component of a
global completed class is a rational squareclass, proving
the asserted nonexistence of such a global class.

For a number field $L$, choose an $E_j$ containing its
Galois closure. Every sufficiently late $q_j$ splits
completely in $L$. Thus at every place $w\mid q_j$ the
image $A_{L,w}$ again has valuation one. The same argument
proves the failure of $P^1(L,\mu_n)$ for every $n>1$, and
excludes a global completed class over $L$.

The inequality $q_j>2^j$ gives
$\sum_jq_j^{-1}<\infty$. Dividing the bounded sum
$\sum_jq_j^{-s}$, for $s>1$, by
$\log(1/(s-1))$ proves that this set has Dirichlet density
zero. At $q=q_j>2$, equation (D.29.4) also gives
\[
              x_q\equiv2\pmod q,
              \qquad 1-x_q\equiv-1\pmod q.
\]
Both original tripod coordinates are therefore units at
precisely the places witnessing the new-coordinate
obstruction. This does not impose their global completed
Kummer conditions.

\emph{All finite-character tests.}
Fix $L,n$ and $\chi$. The kernel of the continuous
character $\chi:G_L\to\Z/n$ cuts out a finite cyclic
extension $K/L$. Choose a finite Galois extension
$N/\Q$ containing $K$ and all its conjugates. There is an
index $J$ such that every $q_j$ for $j\ge J$ splits
completely in $N$. At every $w\mid q_j$ for these indices,
the local character $\chi_w$ is zero.

Let $T$ contain the archimedean places of $L$, the finite
places where $\chi$ ramifies, and all places above the
finitely many $q_j$ with $j<J$. Outside $T$, either $w$
lies above a tail prime, where $\chi_w=0$, or $A_{L,w}$
is a unit and $\chi_w$ is unramified. In the latter case
their cup invariant is zero: the local Artin map for an
unramified extension depends only on the valuation and
annihilates units
\cite[Chapter III, Proposition 3.6 and Lemma 3.7,
pp.~109--110]{D_MilneCFT2020}.
The identical assertion outside $T$ holds for every finite
product $P_r$. This proves that the proposed sum is
well-defined with finite support, even though its first
input is outside the restricted product.

For every finite $w\in T$, the elements $P_r$ converge
to $A_{L,w}$ in $L_w^\times$. The subgroup
$(L_w^\times)^n$ is open. For example Hensel's lemma at
$Z=1$ applied to $Z^n-u$ gives an $n$th root whenever
$v_w(u-1)>2v_w(n)$. Since $T$ is finite, a single sufficiently
large $r$ gives
\[
 \kappa_n(P_r)=\kappa_n(A_{L,w})
                      \quad\text{for every finite }w\in T.
\]
At real places both $P_r>0$ and the chosen value $1$ have
trivial Kummer class, including for even $n$; at complex
places every Kummer class is trivial. We conclude that
\[
\begin{split}
 \sum_w\operatorname{inv}_w
       \bigl(\kappa_n(A_{L,w})\smile\chi_w\bigr)
 &=\sum_w\operatorname{inv}_w
       \bigl(\kappa_n(P_r)\smile\chi_w\bigr)\\
 &=0.
\end{split}
\]
The last equality is global Brauer reciprocity for the
one global class
$\kappa_n(P_r)\smile\chi\in\operatorname{Br}(L)[n]$,
including the real invariants
\cite[Chapter VIII, Theorem 4.2, p.~241]{D_MilneCFT2020}.
The set $T$ depends on $\chi$. This is no assertion that
the first Kummer input is unramified outside a finite set;
that assertion was explicitly disproved.

\emph{Finite-base identities.}
For $L/K$ and $w\mid v$, restriction sends $A_{K,v}$
to $A_{L,w}$, since both are images of the same rational
local limit. The local norm of an element from $K_v$ is
its $[L_w:K_v]$th power. Multiplying over $w\mid v$ yields
\[
 \prod_{w\mid v}N_{L_w/K_v}(A_{L,w})
                         =A_{K,v}^{[L:K]}.
\]
This also holds at the chosen archimedean values and
passes to Kummer completion. It is the usual
norm-of-restriction identity; no stronger coherence is
being assumed.

\emph{Finite certificate and scope.}
For the specified prefix the exact arithmetic is
\begin{center}
\small
\begin{tabular}{rrrrrr}
$j$&$p_j$&$R_j$&$M_j$&$q_j$&$(q_j-1)/M_j$\\
\hline
1&2&2&4&5&1\\
2&3&2&36&37&1\\
3&5&1&1080&2161&2\\
4&7&1&11612160&58060801&5\\
5&11&1&1916006400000&17244057600001&9
\end{tabular}
\end{center}
The program \path{code/D_unrestricted_reciprocity_checks.py}
uses trial division through the integer square root for
every primality decision, verifies the cusp-avoidance
congruences, and checks $45$ quadratic Hilbert reciprocity
products for the finite products. Its exact output is
\path{code/D_unrestricted_reciprocity_checks.json}.
These checks certify the prefix and finite identities; the
proof above establishes the infinite and all-character
assertions.

Poitou--Tate applies to a restricted product. A tuple
outside its domain cannot be placed in its kernel solely
because every individually meaningful character sum
vanishes. Our construction preserves that distinction after
all finite constant extensions. It is not claimed to
satisfy the original global $t,1-t$ conditions, the
conditions on all accessible finite \'etale neighbourhoods,
or any section hypothesis. Accordingly it refutes only
the proposed replacement of admissibility by the tested
reciprocity conditions. The continuous-section obligations
D1--D4 remain unresolved.
\end{proof}
